\section{A coefficient-below-six obstruction and finite selected-place hulls}
\label{sec:conductor-product-hull}

We next close two proposed shortcuts.  On the Frey route, optimization over
an entire rational isogeny class cannot lower the coefficient below six.  On
the local IUT route, projected membership does suffice after the actual
finite product-order hull, but that finite statement does not determine a
global Arakelov degree.

\subsection{An exact conductor and an isogeny-class obstruction}

For $n\geq1$, put $c=1792n+2$ and
\[
 E_c:y^2=x(x-1)(x+c-1)
     =x^3+(c-2)x^2+(1-c)x.
\]
The preceding sections determine the four vertices of its rational isogeny
class and prove that the minimum absolute logarithmic Weil height in that
class is
\begin{equation}\label{eq:new-entire-min-height}
                     3\log(c^2-16c+16).
\end{equation}

\begin{proposition}[Exact local and global conductor]
\label{prop:new-frey-conductor}
The displayed equation is minimal at every prime.  At $2$ it has Kodaira
symbol III and conductor exponent five.  Every odd prime
$s\mid c(c-1)$ is multiplicative of type
$I_{2v_s(c(c-1))}$ and has conductor exponent one.  Hence every elliptic
curve $F/\Q$ rationally isogenous to $E_c$ has
\begin{equation}\label{eq:new-frey-conductor}
              N(F)=16\rad(c(c-1)).
\end{equation}
If $D=(c-1)(c/2)^4$ is the reduced denominator of the height-minimizing
$j$-invariant, then
\begin{equation}\label{eq:new-frey-den-conductor}
                         N(F)=32\rad(D).
\end{equation}
\end{proposition}

\begin{proof}
Write $A=c-2$ and $B=1-c$.  Then
\[
 v_2(A)\geq8,\qquad v_2(B)=0,\qquad B\equiv-1\pmod4.
\]
The corresponding row of the complete two-torsion model table in
Mulholland~\cite[Theorem 2.1 and Table 2.1]{Mulholland2006}, which
specializes Papadopoulos's classification, gives type III, conductor
exponent five, and minimality at two.  The invariants are
\[
 c_4=16(c^2-c+1),\qquad \Delta=16(c-1)^2c^2.
\]
If an odd prime $s$ divides $c$ or $c-1$, then $c^2-c+1\equiv1\pmod s$.
Thus $c_4$ is a unit and the model is minimal multiplicative, with the
asserted type and conductor exponent.  All other odd primes are good.  This
proves \eqref{eq:new-frey-conductor} for $E_c$.  Rational isogenies identify
the rational Tate modules and hence preserve every local Artin conductor,
so the same formula holds throughout the class.

Writing $u=c/2$, one has $D=(2u-1)u^4$ and
$\gcd(u,2u-1)=1$.  Therefore
\[
 \rad(c(c-1))=\rad(2u(2u-1))=2\rad(D),
\]
which gives \eqref{eq:new-frey-den-conductor}.
\end{proof}

The congruence modulus permits an infinite power subfamily.  Put
\[
 c_k=2\cdot897^k,\qquad
 n_k=\frac{897^k-1}{896}=1+897+\cdots+897^{k-1}.
\]
Then $c_k=1792n_k+2$ and, since $897=3\cdot13\cdot23$,
\begin{equation}\label{eq:new-power-conductor}
 N(E_{c_k})=28704\rad(2\cdot897^k-1)<28704c_k.
\end{equation}

\begin{theorem}[Uniform failure of every coefficient below six]
\label{thm:new-frey-six-obstruction}
For every $0\leq\theta<6$ and $C\in\R$, all sufficiently large $k$ satisfy
\[
 h(j(F))>\theta\log N(F)+C
\]
simultaneously for every $F/\Q$ rationally isogenous to $E_{c_k}$.
\end{theorem}

\begin{proof}
For $c_k>32$, \eqref{eq:new-entire-min-height} gives
\[
 h(j(F))>6\log c_k-3\log2.
\]
Equation~\eqref{eq:new-power-conductor} and $\theta\geq0$ give
\[
 \theta\log N(F)<\theta\log c_k+\theta\log28704.
\]
Their difference is greater than
\[
 (6-\theta)\log c_k-3\log2-\theta\log28704,
\]
which tends to infinity because $c_k=2\cdot897^k$ and $\theta<6$.
\end{proof}

This rules out every subcritical coefficient obtained merely by changing
the representative inside the rational isogeny class.  It leaves the
coefficient-six and all-epsilon estimate
$(6+\eps)\log N+C_\eps$ untouched; in particular, it does not prove that
six is an attainable or optimal coefficient.  On this family the
remaining input is a lower bound for
$\rad(2\cdot897^k-1)$; the elementary upper bound in
\eqref{eq:new-power-conductor} cannot supply it.

\subsection{Finite products of selected-place components}

Fix a rational prime and one label.  After decomposing the finite-etale
tensor packet into field factors, write
\[
                         T=\prod_{w\in W}T_w,
             \qquad B=\prod_{w\in W}B_w,
\]
where $W$ is the finite set of ordered selected-place words and $B_w$ is the
product of the integer rings in the field factors of $T_w$.

\begin{proposition}[Exact product-order span]\label{prop:new-product-span}
For every subset $S\subset T$,
\begin{equation}\label{eq:new-product-span}
       \Span_B(S)=\prod_{w\in W}\Span_{B_w}(\pi_wS).
\end{equation}
The same identity holds after topological closure.
\end{proposition}

\begin{proof}
Projection of a $B$-linear combination gives the forward inclusion.  For
the reverse inclusion, let $e_w\in B$ be the central idempotent supported on
the $w$-component.  If
$y_w=\sum_i b_i\pi_w(s_i)$, then
$\sum_i(e_wb_i)s_i$ belongs to $\Span_B(S)$, has component $y_w$ at $w$,
and has every other component zero.  Sum this construction over the finite
set $W$.  Closure commutes with a finite product.
\end{proof}

Suppose now that $S$ satisfies the first-branch hypotheses of
IUT~III, Remark 3.9.5(i): it is relatively compact and contains a relatively
compact finite-log-volume subset.  The source holomorphic hull is the
smallest $\lambda B$ containing $S$, with every field component of $\lambda$
nonzero.  Its definition is componentwise, and therefore
\begin{equation}\label{eq:new-product-hull}
       \operatorname{Hull}_B(S)
       =\prod_{w\in W}\operatorname{Hull}_{B_w}(\pi_wS).
\end{equation}
The compactness hypothesis is essential: in the other branch of the source
definition, the hull is the entire ambient space.

If full-rank product fractional ideals $P_w$ satisfy
\[
 P_w\subset\overline{\Span_{B_w}(\pi_wS)}
 \qquad(w\in W),
\]
then \eqref{eq:new-product-span} gives
\[
       \prod_wP_w\subset\operatorname{Hull}_B(S).
\]
For an ordered word $w=(v_0,\ldots,v_j)$, set
\[
 h_v=[(F_{\rm mod})_v:\Q_p],\qquad
 \omega_w=\frac{\prod_i h_{v_i}}
                 {[F_{\rm mod}:\Q]^{j+1}}.
\]
IUT~IV, Remark 1.7.1 gives precisely these normalized weights.  They sum to
one because
\[
 \sum_w\prod_i h_{v_i}
 =\left(\sum_{v\mid p}h_v\right)^{j+1}
 =[F_{\rm mod}:\Q]^{j+1}.
\]
Monotonicity of Haar measure now yields the finite weighted lower bound
\begin{equation}\label{eq:new-weighted-hull}
 \sum_w\omega_w V_w(\pi_w\operatorname{Hull}(S))
 \geq\sum_w\omega_w V_w(P_w),
\end{equation}
where $V_w=D_w^{-1}\log\mu_w$ and
$D_w=\dim_{\Q_p}T_w$.  No single raw member must realize all word
projections: the central idempotents combine witnesses only after taking the
$B$-span.

Equation~\eqref{eq:new-weighted-hull} is a finite-place result.  It does not
specify the archimedean metric or prove that the finite lattices localize one
global projective module.  Indeed, keep every finite lattice of
$L=\Z\subset\Q$ fixed and set
$\lVert x\rVert_{\infty,t}=e^{-t}|x|$.  All finite projections, hulls, and
Haar volumes are independent of $t$, while
$\widehat\deg(\overline L_t)=t$.  Thus finite projected membership alone
cannot determine a global Arakelov inequality.  The remaining source problem
is to construct the same globally metrized marked family and compare it,
through the actual Ind3 and horizontal links, with the fixed $q$-pilot.
